\documentclass[UTF8,11pt]{ctexart}
\usepackage[a4paper,margin=24mm]{geometry}
\usepackage{amsmath,amssymb,amsthm,mathtools,mathrsfs}
\usepackage[all]{xy}
\usepackage{xurl,hyperref,enumitem}
\usepackage{tikz}
\usetikzlibrary{arrows.meta,cd,decorations.markings,calc,patterns}
\hypersetup{hidelinks,breaklinks=true}
\setlength{\emergencystretch}{3em}
\allowdisplaybreaks
\newtheorem*{theorem}{Theorem}
\newtheorem*{thm}{Theorem}
\newtheorem*{lemma}{Lemma}
\newtheorem*{proposition}{Proposition}
\newtheorem*{prop}{Proposition}
\newtheorem*{corollary}{Corollary}
\newtheorem*{cor}{Corollary}
\newtheorem*{claim}{Claim}
\theoremstyle{definition}
\newtheorem*{definition}{Definition}
\newtheorem*{defn}{Definition}
\newtheorem*{remark}{Remark}
\newtheorem*{example}{Example}
\newcommand{\R}{\mathbb R}
\newcommand{\C}{\mathbb C}
\newcommand{\Z}{\mathbb Z}
\newcommand{\Q}{\mathbb Q}
\newcommand{\N}{\mathbb N}
\newcommand{\F}{\mathbb F}
\newcommand{\D}{\mathbb D}
\newcommand{\bB}{\mathbb B}
\newcommand{\bC}{\mathbb C}
\newcommand{\bR}{\mathbb R}
\newcommand{\bZ}{\mathbb Z}
\newcommand{\bN}{\mathbb N}
\newcommand{\bQ}{\mathbb Q}
\newcommand{\bD}{\mathbb D}
\newcommand{\bF}{\mathbb F}
\newcommand{\CP}{\mathbb{CP}}
\newcommand{\RP}{\mathbb{RP}}
\newcommand{\sO}{\mathscr O}
\newcommand{\sI}{\mathscr I}
\newcommand{\sF}{\mathscr F}
\newcommand{\sG}{\mathscr G}
\newcommand{\sH}{\mathscr H}
\newcommand{\sR}{\mathscr R}
\newcommand{\sM}{\mathscr M}
\newcommand{\sV}{\mathscr V}
\newcommand{\sU}{\mathscr U}
\DeclareMathOperator{\Hom}{Hom}
\DeclareMathOperator{\Ext}{Ext}
\DeclareMathOperator{\Tor}{Tor}
\DeclareMathOperator{\Spec}{Spec}
\DeclareMathOperator{\Proj}{Proj}
\DeclareMathOperator{\supp}{supp}
\DeclareMathOperator{\im}{im}
\DeclareMathOperator{\id}{id}
\DeclareMathOperator{\Tr}{Tr}
\DeclareMathOperator{\tr}{tr}
\DeclareMathOperator{\rank}{rank}
\DeclareMathOperator{\ord}{ord}
\DeclareMathOperator{\dist}{dist}
\DeclareMathOperator{\End}{End}
\DeclareMathOperator{\Aut}{Aut}
\DeclareMathOperator{\Gal}{Gal}
\DeclareMathOperator{\vol}{vol}
\DeclareMathOperator{\Cl}{Cl}
\DeclareMathOperator{\coker}{coker}
\DeclareMathOperator{\codim}{codim}
\DeclareMathOperator{\divergence}{div}
\DeclareMathOperator{\Ric}{Ric}
\DeclareMathOperator{\grad}{grad}
\DeclareMathOperator{\Hess}{Hess}
\newcommand{\abs}[1]{\left\lvert#1\right\rvert}
\newcommand{\sabs}[1]{\lvert#1\rvert}
\newcommand{\babs}[1]{\bigl\lvert#1\bigr\rvert}
\newcommand{\Babs}[1]{\Bigl\lvert#1\Bigr\rvert}
\newcommand{\bbabs}[1]{\biggl\lvert#1\biggr\rvert}
\newcommand{\BBabs}[1]{\Biggl\lvert#1\Biggr\rvert}
\newcommand{\norm}[1]{\left\lVert#1\right\rVert}
\newcommand{\snorm}[1]{\lVert#1\rVert}
\newcommand{\bnorm}[1]{\bigl\lVert#1\bigr\rVert}
\newcommand{\Bnorm}[1]{\Bigl\lVert#1\Bigr\rVert}
\newcommand{\bbnorm}[1]{\biggl\lVert#1\biggr\rVert}
\newcommand{\BBnorm}[1]{\Biggl\lVert#1\Biggr\rVert}
\newcommand{\widebar}[1]{\overline{#1}}
\newcommand{\nicefrac}[2]{\frac{#1}{#2}}
\newcommand{\linnprod}[2]{\langle#1,#2\rangle}
\newcommand{\myindex}[1]{#1}
\newcommand{\glsadd}[1]{}
\newcommand{\avoidbreak}{}
\newcommand{\sourceref}[1]{\textnormal{[原文：\nolinkurl{#1}]}}
\newcommand{\thmref}[1]{\sourceref{#1}}
\newcommand{\lemmaref}[1]{\sourceref{#1}}
\newcommand{\propref}[1]{\sourceref{#1}}
\newcommand{\corref}[1]{\sourceref{#1}}
\newcommand{\defnref}[1]{\sourceref{#1}}
\newcommand{\exampleref}[1]{\sourceref{#1}}
\newcommand{\exerciseref}[1]{\sourceref{#1}}
\newcommand{\figureref}[1]{\sourceref{#1}}
\newcommand{\sectionref}[1]{\sourceref{#1}}
\newcommand{\appendixref}[1]{\sourceref{#1}}
\newcommand{\stacksref}[2]{\href{https://stacks.math.columbia.edu/tag/#1}{#2 [#1]}}


\newcommand{\E}{\mathbb E}
\newcommand{\Prob}{\mathbb P}
\newcommand{\Reff}{\mathcal R_{\mathrm{eff}}}
\newcommand{\eps}{\varepsilon}
\DeclareMathOperator{\diam}{diam}
\DeclareMathOperator{\Var}{Var}
\DeclareMathOperator{\Crit}{Crit}
\newcommand{\dd}{\,\mathrm d}


\begin{document}
\section*{081\quad 强椭圆主符号到Sobolev能量的Garding不等式}
\subsection*{来源与整理范围}
Benoit Dionne, \emph{Partial Differential Equations}, \texttt{elliptic.tex}，Weakly Coercive Bilinear Mappings节，标签\texttt{ell\_garding4, ell\_garding2, ell\_garding3, ell\_garding1}，blob ad9fcaea78e33f90532a08f9bf5801832fa0076c。
\par\url{https://github.com/BenoitDionne/PDE/blob/PDE/elliptic.tex}
\par\textbf{恢复方式（整理者说明）：}依据人类原文重新整理以补齐缺失题号；并非旧题证文件的逐字节恢复；2026-09-28。
\subsection*{作者命题与人类证明}

\paragraph{章内记号（据作者定义整理）。}
$D^\alpha$为分布导数，$H^k=H^{k,2}$，$\|u\|_k^2=\sum_{|\alpha|\leq k}\|D^\alpha u\|_2^2$，
$H^k_0(\Omega)$为$C_c^\infty(\Omega)$的闭包。
内积$\langle v,u\rangle_2=\int_\Omega v\overline u$线性于第一变量。
弱强制性指$\operatorname{Re}B(u,u)\geq C\|u\|_k^2-\lambda\|u\|_2^2$，其中$C>0$。
\begin{theorem}[G\aa rding's inequality; ell\_garding1]
Let $\Omega\subset\mathbb R^n$ be bounded with $C^2$ boundary, and let
$k\geq1$ be an integer. Suppose
\[
Lu=\sum_{|\alpha|,|\beta|\leq k}(-1)^{|\alpha|}
D^\alpha\bigl(b_{\alpha\beta}(x)D^\beta u\bigr)
\]
is strongly elliptic on $\overline\Omega$:
\[
\operatorname{Re}\sum_{|\alpha|=|\beta|=k}b_{\alpha\beta}(x)\xi^\alpha\xi^\beta
\geq\theta|\xi|^{2k}\quad(x\in\overline\Omega,\ \xi\in\mathbb R^n),\quad\theta>0.
\]
Assume $b_{\alpha\beta}\in C(\overline\Omega)$ for $|\alpha|=|\beta|=k$
and $b_{\alpha\beta}\in L^\infty(\Omega)$ for $|\alpha|+|\beta|<2k$.
Then
\[
B(v,u)=\sum_{|\alpha|,|\beta|\leq k}
\langle D^\alpha v,b_{\alpha\beta}D^\beta u\rangle_2
\]
is weakly coercive on $H^k_0(\Omega)$.
\end{theorem}
\paragraph{支撑论证1：作者的插值估计（ell\_garding4）。}
For $r<s<t$ and $\epsilon>0$ there exists $C$ such that
\[
\|f\|_{s,\rho}^2\leq\epsilon\|f\|_{t,\rho}^2+C\|f\|_{r,\rho}^2,
\qquad \|f\|_{s,\rho}^2=\int_{\mathbb R^n}(1+|\xi|^2)^s|\widehat f(\xi)|^2\,d\xi.
\]
\begin{proof}
Choose $R$ so large that $(1+\rho^2)^s\leq\epsilon(1+\rho^2)^t$ for
$\rho>R$. On $|\xi|\leq R$, put $C=(1+R^2)^{s-r}$, so
$(1+|\xi|^2)^s\leq C(1+|\xi|^2)^r$.
In both regions,
\[
(1+|\xi|^2)^s\leq\epsilon(1+|\xi|^2)^t+C(1+|\xi|^2)^r.
\]
Multiply by $|\widehat f|^2$ and integrate. For nonnegative integer
orders, equivalence of the Fourier and derivative Sobolev norms gives
the same assertion for $\|\cdot\|_s$, by first choosing a suitably
smaller $\epsilon$ and adjusting $C$.
\end{proof}
\paragraph{支撑论证2：常系数主部（ell\_garding2）。}
If the strongly elliptic principal coefficients $b_{\alpha\beta}$ are
constant, its form $B_0$ is coercive on $H^k_0(\Omega)$.
\begin{proof}
Extend $u\in H^k_0(\Omega)$ by zero to $\underline u\in H^k(\mathbb R^n)$.
Its distributional derivatives are the zero extensions of those of $u$.
Plancherel gives
\begin{align*}
B_0(u,u)
&=\sum_{|\alpha|=|\beta|=k}\overline{b_{\alpha\beta}}
  \int_{\mathbb R^n}D^\alpha\underline u\,
              \overline{D^\beta\underline u}\,dx\\
&=\int_{\mathbb R^n}
 \left(\sum_{|\alpha|=|\beta|=k}\overline{b_{\alpha\beta}}\xi^\alpha\xi^\beta\right)
 |\widehat{\underline u}(\xi)|^2\,d\xi.
\end{align*}
Taking real parts and using ellipticity,
\begin{align*}
\operatorname{Re}B_0(u,u)
&\geq\theta_1\int |\xi|^{2k}|\widehat{\underline u}(\xi)|^2\,d\xi\\
&\geq\theta_2\sum_{|\alpha|=k}\int
 |\xi^\alpha\widehat{\underline u}(\xi)|^2\,d\xi
=\theta_2\sum_{|\alpha|=k}\|D^\alpha\underline u\|_2^2
\geq\theta_3\|\underline u\|_k^2.
\end{align*}
The last inequality is Poincar\'e's inequality for zero boundary values.
The preceding comparison follows by expanding $(\sum_j\xi_j^2)^k$;
the relevant positive constants are independent of $u$.
Since all these zero extensions vanish off $\Omega$, the last norm
is $\|u\|_{k,\Omega}$, proving coercivity.
\end{proof}
\paragraph{支撑论证3：变系数主部（ell\_garding3）。}
For continuous strongly elliptic principal coefficients, the principal
form $B_1$ is weakly coercive.
\begin{proof}
\textbf{(i) Local coercivity.}
Freeze the coefficients at $z\in\overline\Omega$. The preceding
constant-coefficient bound has a coercivity constant $K>0$ independent
of $z$, by uniform ellipticity. Uniform continuity on $\overline\Omega$
provides $\delta>0$ such that
\[
\sum_{|\alpha|=|\beta|=k}|b_{\alpha\beta}(x)-b_{\alpha\beta}(z)|<K/2
\quad\hbox{when }|x-z|<\delta.
\]
If $u$ is supported in $B_\delta(z)\cap\Omega$, Cauchy--Schwarz gives
\begin{align*}
|B_1(u,u)-B_{1,z}(u,u)|
&\leq\sum_{|\alpha|=|\beta|=k}
 \|b_{\alpha\beta}-b_{\alpha\beta}(z)\|_{L^\infty(\supp u)}
 \|D^\alpha u\|_2\|D^\beta u\|_2\\
&\leq(K/2)\|u\|_k^2.
\end{align*}
Consequently $\operatorname{Re}B_1(u,u)\geq(K/2)\|u\|_k^2$ for such $u$.

\textbf{(ii) Localization and the lower-order terms.}
Choose finitely many such balls covering $\overline\Omega$, and a smooth
partition of unity $\phi_j$ subordinate to them. Set
\[
\psi_j=\phi_j\left(\sum_\ell|\phi_\ell|^2\right)^{-1/2},
\qquad \sum_j|\psi_j|^2=1.
\]
The functions may be chosen real. Writing all sums over
$|\alpha|=|\beta|=k$, we have $B_1(u,u)=Q_1(u)+Q_2(u)$, where
\[
Q_1(u)=\sum_j\sum_{\alpha,\beta}
 \langle D^\alpha(\psi_j u),b_{\alpha\beta}D^\beta(\psi_j u)\rangle_2
\]
and
\begin{align*}
Q_2(u)={}&\sum_{j,\alpha,\beta}
\langle D^\alpha(\psi_j u),
 b_{\alpha\beta}(\psi_jD^\beta u-D^\beta(\psi_j u))\rangle_2\\
&+\sum_{j,\alpha,\beta}
\langle\psi_jD^\alpha u-D^\alpha(\psi_j u),
 b_{\alpha\beta}\psi_jD^\beta u\rangle_2.
\end{align*}
Part (i) gives $\operatorname{Re}Q_1(u)\geq C_1\sum_j\|\psi_ju\|_k^2$.
Expand the derivatives in this latter expression. The top-order terms
without derivatives on $\psi_j$ sum to
\[
\sum_{j,|\alpha|=k}\|\psi_jD^\alpha u\|_2^2
=\sum_{|\alpha|=k}\|D^\alpha u\|_2^2\geq C_2\|u\|_k^2.
\]
All other terms have the form
\[
\int_\Omega A_{\alpha\beta}(x)D^\alpha u\,
\overline{D^\beta u}\,dx,
\qquad |\alpha|,|\beta|\leq k,\quad |\alpha|+|\beta|<2k.
\]
Their coefficients are finite sums of products of derivatives of the
smooth $\psi_j$ and are bounded on $\overline\Omega$.
Cauchy--Schwarz therefore bounds their total absolute value by
$C_3\|u\|_k\|u\|_{k-1}$. We obtain
\[
\operatorname{Re}Q_1(u)\geq C_1C_2\|u\|_k^2-C_1C_3\|u\|_k\|u\|_{k-1}.
\]
In $Q_2$, each commutator
$D^\alpha(\psi_ju)-\psi_jD^\alpha u$ contains only derivatives of $u$
of order at most $k-1$. The same calculation yields
$|Q_2(u)|\leq C_4\|u\|_k\|u\|_{k-1}$.
Thus, with $C'=C_1C_3+C_4$,
\[
\operatorname{Re}B_1(u,u)\geq C_1C_2\|u\|_k^2-C'\|u\|_k\|u\|_{k-1}
\geq\frac{C_1C_2}{2}\|u\|_k^2-
 \frac{(C')^2}{2C_1C_2}\|u\|_{k-1}^2.
\]
Apply the interpolation estimate to the zero extension, with orders
$0,k-1,k$ and $\epsilon=(C_1C_2)^2/(2(C')^2)$, to absorb half the
remaining top-order term. This gives
\[
\operatorname{Re}B_1(u,u)\geq\frac{C_1C_2}{4}\|u\|_k^2-
 \frac{(C')^2C_5}{2C_1C_2}\|u\|_2^2.
\]
If a displayed error constant vanishes, its term is simply absent.
For $k=1$ the lower-order norm is already the $L^2$ norm.
\end{proof}
\begin{proof}[Proof of G\aa rding's inequality]
Write $B=B_1+B_2$, with $B_1$ the principal part and $B_2$ the sum over
$|\alpha|+|\beta|<2k$. The preceding result gives
\[
\operatorname{Re}B_1(u,u)\geq C_1\|u\|_k^2-\lambda_1\|u\|_2^2.
\]
Since all lower-order coefficients are bounded, Cauchy--Schwarz gives
\begin{align*}
|B_2(u,u)|
&\leq\sum_{\substack{|\alpha|,|\beta|\leq k\\|\alpha|+|\beta|<2k}}
 \|b_{\alpha\beta}\|_\infty\|D^\alpha u\|_2\|D^\beta u\|_2\\
&\leq C_2\|u\|_k\|u\|_{k-1}
\leq\frac{C_1}{2}\|u\|_k^2+\frac{C_2^2}{2C_1}\|u\|_{k-1}^2.
\end{align*}
Apply interpolation to the zero extension with
$\epsilon=C_1^2/(2C_2^2)$, obtaining
\[
\frac{C_2^2}{2C_1}\|u\|_{k-1}^2
\leq\frac{C_1}{4}\|u\|_k^2+\frac{C_2^2C_3}{2C_1}\|u\|_2^2.
\]
(The zero-error and $k=1$ cases are read as above.) Hence
\[
\operatorname{Re}B(u,u)\geq\frac{C_1}{4}\|u\|_k^2-
\left(\lambda_1+\frac{C_2^2C_3}{2C_1}\right)\|u\|_2^2,
\]
which is the required weak coercivity.
\end{proof}

\subsection*{整理说明与先修范围（非作者正文）}
收录作者的Fourier插值、常系数、系数冻结/单位分解、一般低阶项四段论证，重复的有限多重指标估计合并排版；各交换子和误差吸收步骤保留。标准先修是零延拓、Plancherel、Sobolev范数等价及Poincar\'e不等式。原文局部界漏印平方、平方单位分解的归一化指数写为$-1$、$Q_2$只写单侧界、末尾$\lambda$多写一个分母因子，此处分别按同段计算改为平方、$-1/2$、绝对值界及显示式中的系数；插值范数换算统一用可变常数；$k=1$和零误差常数的读法单独注明。数学路线未改。难点是将主符号逐点正性转成全局能量下界，同时控制局部化造成的低阶误差。
\subsection*{可修改的简短标注（非作者正文）}
$c$：高阶椭圆算子的弱强制能量估计

$a$：Fourier变换；主符号；冻结系数；平方单位分解；微分交换子；Sobolev插值

\texttt{cross\_theory\_status = yes}。常系数段把微分能量转为频率多项式正性；变系数段再用平方单位分解把局部信息拼回全局，并将不保持主阶的误差降至可吸收的低阶范数。

全部标注为非穷尽、可修改初稿。
\subsection*{来源许可与编辑归属}
Benoit Dionne，\emph{Partial Differential Equations}，University of Ottawa，README署年2024，CC BY-NC-SA 4.0：\url{https://github.com/BenoitDionne/PDE/blob/PDE/README.md}；\url{https://creativecommons.org/licenses/by-nc-sa/4.0/}。本题编辑版本遵循同一许可。
ChatGPT 加入题号、中文来源说明、整理范围和初稿标注；编辑说明与人类证明分列。
\end{document}
